% secV7_goal.tex -- Section 4. CONTAINER (27 Oct 2026): text to be written by the author.
% Word budget: 2,000 words (JPART count: running text, footnotes and citations as printed).
% Outline, numbers and sources: referee_EE_2026-10/ESQUELETO.md, section 4. The archived text
% (_arquivo_texto_IA_2026-10-27/) is not a source.
% \label{sec:pandemic} is defined in this section, on the subsection on the post-2019 period and the
% official review; the note of Table A13 (V3w_indet_full) refers to it.

\section{What Goal~5 leaves open}\label{sec:goal}

Goal~5 of S\~ao Paulo's Municipal Education Plan requires the provision of childcare to at least 75\% of children aged 0--3 years and 11 months or to 100\% of registered demand, whichever is greater \citep{saopaulo2015pme}. The law leaves the operational conventions needed to assess compliance unspecified. Across our measurement grid, Goal~5 compliance is specification-dependent in 2019--2022, whereas no specification in that grid indicates compliance in any evaluable year in 2010--2018; before September 2019, only the published waiting-list series is available.

\subsection{Operationalizing the two requirements}

Let $P_t$ denote the estimated number of children below age four in year $t$, $E^A_t$ the number enrolled under the coverage definition used for the population requirement, $E^B_t$ the number enrolled under the administrative definition used for registered demand, and $L_t$ the number of children remaining on the relevant waiting list. The two requirements imply quantities $T^A_t=0.75P_t$ and $T^B_t=E^B_t+L_t$. A natural implementation compares $T^A_t$ and $T^B_t$, assesses $E^A_t\geq 0.75P_t$ when the first is larger, and requires $L_t=0$ when the second is larger. The latter condition is equivalent, for positive enrolment, to $E^B_t/(E^B_t+L_t)=1$. This rule is an analytical reconstruction, not a statutory formula. Because $E^A_t$ and $E^B_t$ may differ in universe or date, their comparison requires another convention absent from the law.

The waiting list presents the first ambiguity. The published series includes centre-specific requests even when families are offered places elsewhere; a narrower series, available only from September 2019, excludes them. These families therefore count as having unmet demand only in the published series. Across 89 published snapshots, the full list is never zero. By contrast, the narrower list is zero in each December observation from 2020 through 2025, while remaining positive in the other 30 observations between September 2019 and August 2026. Whether the queue is empty thus depends on its definition, and the series cover different periods.

Goal~5 does not specify a reference date. A subsequent municipal monitoring ordinance records indicators as of December 31 \citep{sme2016portaria7972}, and the plan's baseline report uses December 31, 2015 \citep{sme2016linhabase}; December is consequently our principal specification. In 2019, the published waiting-list count fell from 75,267 in September to 9,670 in December. We examine September snapshots and annual averages in the appendix.

The enrolment quantities require distinct choices. For the population requirement, the School Census permits a count of children enrolled in the childcare stage across all school networks, or a narrower count covering municipal and publicly partnered providers. Article~7 of the municipal law states that the plan's goals apply to all educational systems in the municipality. We read this as supporting the all-network interpretation for the main analysis, while recognizing that ``systems'' and ``networks'' are not defined as equivalent in the statute; the municipal-and-partnered interpretation remains a sensitivity check. The choice matters: using our principal population series, all-network childcare-stage enrolment reaches 76.8\% of children under four in 2021, whereas the municipal-and-partnered measure reaches only 70.3\%. A classification break prevents a comparable municipal-and-partnered series after 2021. An alternative age-based count, including enrolment at any stage, yields 78.6\% in 2022, against 83.8\% for the childcare-stage measure (Table~\ref{tab:arms}).

For the registered-demand requirement, enrolment may come from the municipal secretariat's December publication, which also reports the waiting list, or from municipal-and-partnered enrolment in the May School Census. The sources differ in date and coverage. In 2021 the secretariat reported 385,100 municipal childcare enrolments in December, exceeding even the 377,862 enrolments across all networks reported by the May School Census. This need not indicate error in either source. Finally, ``whichever is greater'' can mean testing only the larger requirement, as in our main specification, or requiring an empty list regardless. We retain both readings.

\subsection{The denominator is not specified by the law}

The population denominator is a sixth, independent choice. Article~4 of the municipal law refers to the most recent \emph{basic and higher education censuses available on the date of the law's publication} (September 17, 2015); it does not select a demographic census, a population projection, or a procedure for estimating children under four. Our principal series combines civil registration of live births with the 2010 and 2022 Demographic Censuses. For a School Census observation in May of year $t$, we count births from June of $t-4$ through May of $t$ and divide this four-year pool by a scaling factor. The factor is the ratio of the corresponding birth pool to the census population aged 0--3 at the two benchmark years (1.228 in 2010 and 1.278 in 2022), interpolated linearly between them. By construction the resulting series reproduces the two demographic census benchmarks. It preserves cohort variation while assuming a linear scaling factor.

We compare this measure with two alternatives (Table~\ref{tab:armArob}). The first interpolates the 0--3 population directly between the two Demographic Censuses, from 564,210 in 2010 to 464,826 in 2022, without using the intervening birth series. The second interpolates published SEADE projections. Our grid admits two series types: reconstructions calibrated to both Demographic Census benchmarks, and interpolations of independently published population projections. This provenance criterion does not select series by their proximity to a single benchmark. The SEADE interpolation is included because the secretariat cites SEADE estimates when describing coverage, although Goal~5 neither requires nor excludes them. Our interpolation of SEADE's published projections yields 531,819 children in 2022, 14.4\% above the 2022 Demographic Census count of 464,826; this is a discrepancy, not an estimate of SEADE's error. A SEADE series mechanically rescaled to match only the 2022 Census satisfies neither criterion: it gives 498,176 children in 2010, 11.7\% below the Census count of 564,210. We exclude it from the grid, but note that it yields compliance in 2018 under two December published-list specifications. This exclusion is not mandated by law.

These denominators matter near the threshold. With the births-scaled series, all-network childcare-stage enrolment rises from 67.4\% in 2018 to 73.0\% in 2019, 74.1\% in 2020, 76.8\% in 2021, and 83.8\% in 2022 (Table~\ref{tab:arms}). A population 2.6\% lower in 2019 or 1.2\% lower in 2020 would move the first requirement to 75\%; conversely, a population 2.4\% higher in 2021 would move it back to the threshold. For perspective, the post-enumeration evaluation of Brazil's 2022 Census reports a national net coverage error of 8.3\%, with estimates around 10\% for children aged up to four \citep{agenciabrasil2024ppe}. These national rates cannot be applied as corrections for S\~ao Paulo, but underscore why near-threshold classifications are sensitive to measurement.

\subsection{Compliance across specifications}

Table~\ref{tab:indet} evaluates 12 combinations holding constant the December reference count, whatever its printed extraction date, all-network enrolment for the population requirement, and the rule that only the larger of the two requirements binds. The combinations vary the published versus narrower list, the three population series, and the secretariat versus School Census enrolment measure for registered demand. Seven combinations meet the goal in at least one observable year; five meet it in none. The earliest qualifying year is 2019 in four combinations and 2021 in three. None classifies a year before 2019 as compliant within the observations that can be evaluated.

The classifications differ across these combinations. With the published list, births-scaled population, and secretariat enrolment, the goal is not met in any evaluable year; replacing secretariat enrolment with the School Census changes the classification for 2021. Using the linearly interpolated population instead gives compliance in 2019 alone with secretariat enrolment and in 2019--2021 with School Census enrolment. Under the narrower list and births-scaled population, secretariat enrolment yields compliance in 2021--2022, whereas the School Census yields 2021 only. Combining the narrower list with the linear population series gives compliance in 2019--2022 under secretariat enrolment and in 2019--2021 under School Census enrolment. None of the four combinations using the SEADE projection meets the goal. The published-list, linear-population, secretariat-enrolment specification qualifies in 2019 but not later, yielding a nonmonotonic classification.

{\emergencystretch=1em Table~\ref{tab:indetfull} extends this exercise to September observations, annual-average lists, municipal-and-partnered enrolment for the population requirement, and the interpretation that registered demand must always be fully served. The resulting grid contains 144 combinations. The grid enumerates specifications; it is neither a distribution of legally valid interpretations nor a statistical confidence interval \citep{steegen2016increasing,simonsohn2020specification}. Of the 144 combinations, 21 meet the goal in at least one evaluable year and 123 never do; four exhibit a transition from compliance to noncompliance in a later evaluable year. Among combinations that ever qualify, four first do so in 2019, eight in 2020, and nine in 2021. None first qualifies in 2022. In the 2010--2018 period, 72 combinations can be evaluated, all necessarily based on the published list because the narrower series did not yet exist. None of the resulting 648 combination-year observations meets the goal. Evaluable years differ: the narrower list starts in 2019 and the relevant School Census split ends in 2021. A ``never'' entry means never among observed, evaluable years, not a verified failure at all dates.\par}

\subsection{The post-2019 period and the official review}\label{sec:pandemic}

All positive compliance classifications in our specification grid occur in 2019--2022, a period combining enrolment changes, declining birth cohorts, and the disruption of the COVID-19 pandemic. The published September waiting list falls from 75,267 in 2019 to 6,670 in 2020, while secretariat enrolment increases by 20,490. Enrolment plus the published waiting-list count consequently falls by 11.7\%, compared with a 2.8\% decline in the relevant birth pool. The December list falls from 9,670 in 2019 to 540 in 2020; over the same interval, December municipal enrolment increases by 29,179. Births in 2020, 2021, and 2022 are, respectively, 7.2\%, 13.8\%, and 16.7\% below their 2019 level. These movements do not identify changes in preferences, applications, capacity, or pandemic effects.

Without post-2018 individual applications, we cannot distinguish fewer registrations from greater outflows or recording changes. Stock accounting provides a lower bound on first-half entries: increases from December to June imply at least 29,213 entries in 2019, 13,062 in 2020, 8,227 in 2021, and 4,166 in 2022, even before allowing for exits during the interval. Where individual records exist, this bound captures only 24--42\% of first requests. It establishes continuing inflows but identifies neither pandemic effects nor counterfactual compliance. Furthermore, 14 of the 21 specifications that ever qualify do so only in 2020--2021, but none of those 14 can be evaluated in 2022. Their observed patterns cannot be interpreted as evidence of a subsequent reversal.

The municipality's 2023 review nonetheless classified Goal~5 as achieved and reported that the 75\% threshold had been reached as early as 2015 \citep{saopaulo2023conferencia}. Its published indicator was labelled the share of registered childcare demand served, with December as the reference month, rather than an explicit test of the legal maximum of the two requirements. The review did not publish the indicator's formula. We find that $E^B_t/(E^B_t+L_t)$, constructed from December municipal enrolments and the published list, reproduces six of its eight reported values for 2015--2022 to two decimal places. The exceptions are 2019 (96.31\% reported versus 97.28\% reconstructed) and 2022 (99.99\% versus approximately 99.81\%). Published documentation does not explain the discrepancies. Crucially, this indicator neither identifies the binding statutory requirement nor tests compliance with it. The review's classification thus cannot be recovered as compliance with any of the 144 specifications examined here, even though much of its reported series is numerically close to a recognizable administrative ratio.

The availability of these retrospective specifications to the 2023 reviewers is unknown: the publication date of the age-specific 2022 Census data has not been verified. None is a legally definitive indicator. Rather, the same statutory wording supports materially different empirical classifications, and the administrative review does not document the additional conventions necessary to adjudicate among them. A reproducible judgement about Goal~5 therefore requires making these six conventions explicit.
